% Options for packages loaded elsewhere
\PassOptionsToPackage{unicode}{hyperref}
\PassOptionsToPackage{hyphens}{url}
\documentclass[
  ignorenonframetext,
]{beamer}
\newif\ifbibliography
\usepackage{pgfpages}
\setbeamertemplate{caption}[numbered]
\setbeamertemplate{caption label separator}{: }
\setbeamercolor{caption name}{fg=normal text.fg}
\beamertemplatenavigationsymbolsempty
% remove section numbering
\setbeamertemplate{part page}{
  \centering
  \begin{beamercolorbox}[sep=16pt,center]{part title}
    \usebeamerfont{part title}\insertpart\par
  \end{beamercolorbox}
}
\setbeamertemplate{section page}{
  \centering
  \begin{beamercolorbox}[sep=12pt,center]{section title}
    \usebeamerfont{section title}\insertsection\par
  \end{beamercolorbox}
}
\setbeamertemplate{subsection page}{
  \centering
  \begin{beamercolorbox}[sep=8pt,center]{subsection title}
    \usebeamerfont{subsection title}\insertsubsection\par
  \end{beamercolorbox}
}
% Prevent slide breaks in the middle of a paragraph
\widowpenalties 1 10000
\raggedbottom
\AtBeginPart{
  \frame{\partpage}
}
\AtBeginSection{
  \ifbibliography
  \else
    \frame{\sectionpage}
  \fi
}
\AtBeginSubsection{
  \frame{\subsectionpage}
}
\usepackage{iftex}
\ifPDFTeX
  \usepackage[T1]{fontenc}
  \usepackage[utf8]{inputenc}
  \usepackage{textcomp} % provide euro and other symbols
\else % if luatex or xetex
  \usepackage{unicode-math} % this also loads fontspec
  \defaultfontfeatures{Scale=MatchLowercase}
  \defaultfontfeatures[\rmfamily]{Ligatures=TeX,Scale=1}
\fi
\usepackage{lmodern}
\ifPDFTeX\else
  % xetex/luatex font selection
\fi
% Use upquote if available, for straight quotes in verbatim environments
\IfFileExists{upquote.sty}{\usepackage{upquote}}{}
\IfFileExists{microtype.sty}{% use microtype if available
  \usepackage[]{microtype}
  \UseMicrotypeSet[protrusion]{basicmath} % disable protrusion for tt fonts
}{}
\makeatletter
\@ifundefined{KOMAClassName}{% if non-KOMA class
  \IfFileExists{parskip.sty}{%
    \usepackage{parskip}
  }{% else
    \setlength{\parindent}{0pt}
    \setlength{\parskip}{6pt plus 2pt minus 1pt}}
}{% if KOMA class
  \KOMAoptions{parskip=half}}
\makeatother
\usepackage{longtable,booktabs,array}
\newcounter{none} % for unnumbered tables
\usepackage{calc} % for calculating minipage widths
\usepackage{caption}
% Make caption package work with longtable
\makeatletter
\def\fnum@table{\tablename~\thetable}
\makeatother
\setlength{\emergencystretch}{3em} % prevent overfull lines
\providecommand{\tightlist}{%
  \setlength{\itemsep}{0pt}\setlength{\parskip}{0pt}}
\usepackage{bookmark}
\IfFileExists{xurl.sty}{\usepackage{xurl}}{} % add URL line breaks if available
\urlstyle{same}
\hypersetup{
  hidelinks,
  pdfcreator={LaTeX via pandoc}}

\author{\texorpdfstring{}{}}
\date{}

\begin{document}

\section{Exports, parser capabilities, and progressive IR
stages}\label{exports-parser-capabilities-and-progressive-ir-stages}

\begin{frame}[fragile]{Export targets}
\protect\phantomsection\label{export-targets}
The primary export targets are the \textbf{Generalized Notation Notation
(GNN)} canonical Markdown (\texttt{model.gnn.md}) and the equivalent
companion JSON files described in
\texttt{../cogant/docs/export/README.md}. Optional interop targets
(GraphML, Parquet) support analysis in Gephi/yEd and DuckDB, and
optional tensor views for PyTorch Geometric, DGL, or HDF5 can be
selected when downstream graph neural network training pipelines need to
consume the program graph as a relational tensor. Ensure the Python
environment includes optional dependencies for these tensor exports when
those code paths are used.
\end{frame}

\begin{frame}[fragile]{Python AST parser capabilities}
\protect\phantomsection\label{python-ast-parser-capabilities}
The v0.5.x front end relies on
\texttt{cogant.static.parser.PythonASTParser}, which processes Python
source through the standard-library \texttt{ast} module at the CPython
version available in the runtime (3.11+ required, consistent with the
\texttt{requires-python\ =\ "\textgreater{}=3.11"} declared in
\texttt{../cogant/pyproject.toml}). The parser extracts the following
construct categories:

\begin{itemize}
\tightlist
\item
  \textbf{Module-level entities}: module docstrings,
  \texttt{\_\_all\_\_} exports, top-level assignments.
\item
  \textbf{Functions and methods}: \texttt{def} and \texttt{async\ def},
  including signatures with positional, keyword, variadic
  (\texttt{*args}, \texttt{**kwargs}), and positional-only parameters.
  Default values are recorded as constant expressions where statically
  evaluable.
\item
  \textbf{Classes}: class definitions, base classes, metaclasses, and
  the \texttt{\_\_init\_\_} / \texttt{\_\_new\_\_} boundary.
\item
  \textbf{Decorators}: \texttt{@staticmethod}, \texttt{@classmethod},
  \texttt{@property}, \texttt{@dataclass}, and arbitrary user-defined
  decorators. Decorator arguments are captured as attribute metadata.
\item
  \textbf{Type annotations}: PEP 484 / 526 / 604 annotations on function
  parameters, return types, and variable assignments. Generic subscripts
  (\texttt{List{[}int{]}}, \texttt{Dict{[}str,\ Any{]}}) are preserved
  as type strings.
\item
  \textbf{Comprehensions and generators}: list, set, dict comprehensions
  and generator expressions are represented as anonymous FUNCTION nodes
  with DATA\_FLOW edges to their enclosing scope.
\item
  \textbf{Control flow}: \texttt{if}/\texttt{elif}/\texttt{else},
  \texttt{for}/\texttt{while}/\texttt{else},
  \texttt{try}/\texttt{except}/\texttt{finally},
  \texttt{with}/\texttt{async\ with}, and \texttt{match}/\texttt{case}
  (Python 3.10+) are mapped to CONTROLFLOW\_NODE entities.
\item
  \textbf{Imports}: \texttt{import} and \texttt{from\ ...\ import}
  statements produce MODULE\_IMPORT roles with edges to the resolved
  module when discoverable on the file system.
\item
  \textbf{Constants}: module-level and class-level assignments to
  \texttt{Final} or ALL\_CAPS names are classified as CONSTANT nodes.
\end{itemize}

Constructs that require runtime evaluation (for example \texttt{exec},
\texttt{importlib.import\_module}, or dynamic \texttt{\_\_getattr\_\_})
are recorded as EXTERNAL nodes with HEURISTIC provenance and
correspondingly lower confidence.
\end{frame}

\begin{frame}{Progressive IR stages}
\protect\phantomsection\label{progressive-ir-stages}
Processing advances through six intermediate representations, each
adding semantic detail atop its predecessor. Table 3 summarizes what
each stage contributes.

\textbf{Table 3. Progressive IR stages and their contributions.}

{\def\LTcaptype{none} % do not increment counter
\begin{longtable}[]{@{}
  >{\raggedright\arraybackslash}p{(\linewidth - 6\tabcolsep) * \real{0.0986}}
  >{\raggedright\arraybackslash}p{(\linewidth - 6\tabcolsep) * \real{0.1268}}
  >{\raggedright\arraybackslash}p{(\linewidth - 6\tabcolsep) * \real{0.2113}}
  >{\raggedright\arraybackslash}p{(\linewidth - 6\tabcolsep) * \real{0.5634}}@{}}
\toprule\noalign{}
\begin{minipage}[b]{\linewidth}\raggedright
Stage
\end{minipage} & \begin{minipage}[b]{\linewidth}\raggedright
IR name
\end{minipage} & \begin{minipage}[b]{\linewidth}\raggedright
Key additions
\end{minipage} & \begin{minipage}[b]{\linewidth}\raggedright
Typical output size (10K-function repo)
\end{minipage} \\
\midrule\noalign{}
\endhead
1 & Repo IR & Raw entities and relationships per file; deduplication;
merged type info & \textasciitilde15 MB JSON \\
2 & Program Graph IR & Consolidated directed graph \(G=(V,E)\); stable
identifiers; confidence and provenance on every node and edge &
\textasciitilde20 MB JSON \\
3 & Semantic Mapping IR & Translation rules applied; semantic roles
assigned; confidence adjusted by rule engine & \textasciitilde22 MB JSON
(graph + mapping log) \\
4 & State Space IR & Variables, actions, transitions, observations;
dynamic traces integrated where available & \textasciitilde5 MB JSON
(behavioral model) \\
5 & Process Model IR & Higher-level control patterns (request--response,
producer--consumer, state machines) & \textasciitilde2 MB JSON \\
6 & Validation IR & Coverage metrics, confidence distribution, schema
compliance, consistency checks, reproducibility hashes &
\textasciitilde1 MB JSON (report) \\
\bottomrule\noalign{}
\end{longtable}
}

Stages 4 and 5 are \textbf{partial} for many repositories: the
state-space compiler requires either execution traces or sufficient
static structure (for example annotated state machines) to produce
meaningful output. Where dynamic evidence is available, COGANT's
ingestion pipeline follows the established pattern of attaching runtime
observations (coverage, call frequencies, traces) to static program
elements --- dynamic instrumentation frameworks such as Pin
{[}@luk2005pin{]} and invariant detectors such as Daikon
{[}@ernst2007daikon{]} established this general approach of augmenting
static program structure with execution-time evidence. The pipeline
tolerates missing stages gracefully; the Validation IR records which
stages completed and which were skipped.
\end{frame}

\end{document}
